\chapter{Higher-Order Partials, Young's Theorem, the Hessian, and Definiteness}

First-order partials describe slopes; second-order partials describe
\emph{curvature}. In one variable we used $f''$ to distinguish maxima from
minima. In several variables we need a whole \emph{matrix} of second
derivatives (the Hessian) and curvature is captured by its
definiteness, which we test with the leading principal minors (Sylvester's
criterion).

\section{Core Rules}

\begin{rulebox}{Direct second-order partials}
For $f(x,y)$,
\[
  f_{xx} \;=\; \pd{}{x}\!\left(\pd{f}{x}\right) \;=\; \frac{\partial^{2} f}{\partial x^{2}},
  \qquad
  f_{yy} \;=\; \frac{\partial^{2} f}{\partial y^{2}} .
\]
\end{rulebox}

\begin{rulebox}{Cross (mixed) partials}
\[
  f_{xy} \;=\; \pd{}{y}\!\left(\pd{f}{x}\right) \;=\; \frac{\partial^{2} f}{\partial y\,\partial x},
  \qquad
  f_{yx} \;=\; \pd{}{x}\!\left(\pd{f}{y}\right) .
\]
\end{rulebox}

\begin{rulebox}{Young's Theorem (Schwarz / Clairaut)}
If $f$ is $C^{2}$ on an open set (i.e.\ all second partials exist and
are continuous) then the order of differentiation is irrelevant:
\[
  f_{xy} \;=\; f_{yx} .
\]
This is the symmetry of the Hessian. All standard polynomial, exponential,
logarithmic, Cobb--Douglas, and CES forms are $C^{2}$ on their natural
domains.
\end{rulebox}

\begin{rulebox}{Hessian matrix}
For $f:\R^{n}\to\R$,
\[
  \Hess_f(\mathbf{x}) \;=\; \bigl[f_{x_ix_j}(\mathbf{x})\bigr]_{i,j=1}^{n}
  \;=\;
  \begin{pmatrix}
    f_{x_1x_1} & f_{x_1x_2} & \cdots & f_{x_1x_n} \\
    f_{x_2x_1} & f_{x_2x_2} & \cdots & f_{x_2x_n} \\
    \vdots     & \vdots     & \ddots & \vdots     \\
    f_{x_nx_1} & f_{x_nx_2} & \cdots & f_{x_nx_n}
  \end{pmatrix} .
\]
For $C^{2}$ functions $\Hess_f$ is symmetric.
\end{rulebox}

\begin{rulebox}{Definiteness}
A symmetric matrix $H$ is
\begin{itemize}[leftmargin=1.6em,itemsep=0.3ex]
  \item \emph{positive definite (PD)} if $\mathbf{v}^{T}H\mathbf{v}>0$ for
        every $\mathbf{v}\neq\mathbf{0}$ (curvature ``up'' in every
        direction $\Rightarrow$ local minimum;
  \item \emph{negative definite (ND)} if $\mathbf{v}^{T}H\mathbf{v}<0$ for
        every $\mathbf{v}\neq\mathbf{0}$) curvature ``down'' in every
        direction $\Rightarrow$ local maximum;
  \item \emph{indefinite} if $\mathbf{v}^{T}H\mathbf{v}$ takes both
        positive and negative values $\Rightarrow$ saddle point.
\end{itemize}
\end{rulebox}

\begin{rulebox}{Sylvester's criterion ($2\times2$ case)}
Let $H=\begin{pmatrix}a&b\\b&c\end{pmatrix}$ with leading principal minors
$D_1=a$ and $D_2=\det H=ac-b^{2}$. Then
\[
  \begin{aligned}
    D_1>0 \text{ and } D_2>0 &\;\Rightarrow\; H \text{ is PD}, \\
    D_1<0 \text{ and } D_2>0 &\;\Rightarrow\; H \text{ is ND}, \\
    D_2<0 &\;\Rightarrow\; H \text{ is indefinite (saddle)}, \\
    D_2=0 &\;\Rightarrow\; \text{test is inconclusive (semi-definite case)}.
  \end{aligned}
\]
Note: ND requires $D_1<0$ \emph{and} $D_2>0$; the alternating-sign
condition is what one generalises to higher $n$ (signs $-,+,-,+,\dots$ for
ND).
\end{rulebox}

\begin{rulebox}{Classification of critical points}
A point $\mathbf{x}_0$ is a \emph{critical point} of $f$ if
$\grad f(\mathbf{x}_0)=\mathbf{0}$. Then evaluate $\Hess_f(\mathbf{x}_0)$:
\begin{itemize}[leftmargin=1.6em,itemsep=0.3ex]
  \item $H$ PD at $\mathbf{x}_0$ $\Rightarrow$ strict local minimum;
  \item $H$ ND at $\mathbf{x}_0$ $\Rightarrow$ strict local maximum;
  \item $H$ indefinite at $\mathbf{x}_0$ $\Rightarrow$ saddle point;
  \item $D_2=0$ $\Rightarrow$ inconclusive, higher-order analysis
        required.
\end{itemize}
\end{rulebox}

\section{Worked Example}

\begin{worked}{Find and classify every critical point of $f(x,y)=x^{3}-3xy+y^{3}$}
\wstep{Step 1. First-order conditions: both partials zero, simultaneously.}
\[
  f_{x}=3x^{2}-3y=0
  \quad\Longrightarrow\quad y=x^{2},
\]
\[
  f_{y}=-3x+3y^{2}=0
  \quad\Longrightarrow\quad x=y^{2} .
\]
\wtrap{A critical point needs \emph{every} partial to vanish at the
\emph{same} point. Solving each equation separately and pairing the answers
arbitrarily is not the same thing --- the two equations form a system.}

\wstep{Step 2. Solve the system by substitution.}
Substituting $y=x^{2}$ into $x=y^{2}$:
\[
  x=x^{4}
  \;\Longrightarrow\;
  x^{4}-x=0
  \;\Longrightarrow\;
  x\bigl(x^{3}-1\bigr)=0
  \;\Longrightarrow\;
  x=0\ \text{or}\ x=1 .
\]
The corresponding $y=x^{2}$ gives the two critical points
\[
  (0,0)\qquad\text{and}\qquad(1,1).
\]
\wtrap{Do not divide by $x$ to get $x^{3}=1$. That silently discards the
root $x=0$ --- and $(0,0)$ turns out to be the more interesting of the two
points. Factor; never divide by something that may vanish.}

\wstep{Step 3. Build the Hessian once, symbolically.}
\[
  f_{xx}=6x,\qquad f_{yy}=6y,\qquad f_{xy}=f_{yx}=-3 ,
\]
\[
  \Hess f(x,y)=\begin{pmatrix}6x&-3\\-3&6y\end{pmatrix}.
\]
The two mixed partials agree, as Young's theorem guarantees for a function
with continuous second derivatives --- differentiating in $x$ then $y$ gives
the same result as $y$ then $x$. Computing both is a free check, and here
$\frac{\partial}{\partial y}(3x^{2}-3y)=-3$ and
$\frac{\partial}{\partial x}(-3x+3y^{2})=-3$. \checkmark

\wstep{Step 4. Classify $(0,0)$.}
\[
  \Hess f(0,0)=\begin{pmatrix}0&-3\\-3&0\end{pmatrix},
  \qquad
  \det=0-9=-9<0 .
\]
A negative determinant means the two eigenvalues have opposite signs, so the
form is indefinite:
\[
  \boxed{(0,0)\ \text{is a saddle point.}}
\]
No further test is needed, and the value of $f_{xx}$ is irrelevant here.
\wtrap{$\det\Hess<0$ settles the matter immediately --- saddle, whatever
$f_{xx}$ does. It is only when $\det\Hess>0$ that $f_{xx}$ has to be
consulted. Checking $f_{xx}$ first wastes effort and invites the wrong
conclusion.}

\wstep{Step 5. Classify $(1,1)$.}
\[
  \Hess f(1,1)=\begin{pmatrix}6&-3\\-3&6\end{pmatrix},
  \qquad
  \det=36-9=27>0,
  \qquad
  f_{xx}=6>0 .
\]
Positive determinant with a positive $f_{xx}$: positive definite, so
\[
  \boxed{(1,1)\ \text{is a strict local minimum},\qquad f(1,1)=1-3+1=-1 .}
\]

\wstep{Step 6. Verify the saddle directly, because it is cheap.}
Along the line $y=x$: $f(x,x)=2x^{3}-3x^{2}$, whose value near $0$ is
$\approx-3x^{2}<0$. Along $y=-x$: $f(x,-x)=x^{3}+3x^{2}-x^{3}=3x^{2}>0$. So
$f$ takes both signs arbitrarily close to $(0,0)$, where $f=0$ --- a saddle,
confirmed without any matrix.

\wstep{Step 7. The decision rule, complete.}
\begin{center}\footnotesize
\begin{tabular}{lll}
  \hline
  $\det\Hess$ & $f_{xx}$ & conclusion\\
  \hline
  $>0$ & $>0$ & strict local minimum\\
  $>0$ & $<0$ & strict local maximum\\
  $<0$ & --- & saddle point\\
  $=0$ & --- & \textbf{test fails}; investigate directly\\
  \hline
\end{tabular}
\end{center}
\wtrap{$\det\Hess=0$ means the test is silent, not that the point is a
saddle. $f=x^{4}+y^{4}$ (minimum), $f=-x^{4}-y^{4}$ (maximum) and
$f=x^{4}-y^{4}$ (saddle) all have the zero Hessian at the origin. When the
determinant vanishes you must argue directly, as in Step~6.}

\wstep{Step 8. Local is not global.}
Neither point is a global extremum: $f(x,0)=x^{3}$ runs to $\pm\infty$. On
an unbounded domain, second-order conditions can never certify a global
optimum by themselves --- that needs convexity of $f$ on the whole domain,
or a compact feasible set and Weierstrass. Volume~7 makes this the
difference between finding a candidate and having an answer.
\end{worked}

\section{Practice Problems}

\begin{problems}
  \item Compute all second partials of $f(x,y) = 4x^{3} - 5x^{2}y + 7y^{4}$
        and verify Young's theorem.
  \item Construct the Hessian of
        $f(x,y) = 5x^{2} + 6xy + 2y^{2} - 3x + 4y$.
  \item Classify the definiteness of
        $\Hess = \begin{pmatrix}6&2\\2&3\end{pmatrix}$ using Sylvester's
        criterion.
  \item Find and classify the critical points of
        $f(x,y) = x^{2} + 4y^{2} - 6x + 8y + 11$.
  \item Compute all second partials of $f(x,y) = (x+2y)^{3}$ and verify
        $f_{xy} = f_{yx}$.
  \item Construct the Hessian of $f(x,y) = x^{3} + 2xy^{2} - 4y$.
  \item Classify $\Hess = \begin{pmatrix}-4&2\\2&-5\end{pmatrix}$.
  \item Find and classify the critical points of
        $f(x,y) = -x^{2} - 3y^{2} + 8x - 6y + 5$.
  \item Compute all second partials of $f(x,y) = x^{2}\ln(y)$ ($y>0$) and
        verify Young's theorem.
  \item Construct the Hessian of $f(x,y) = \mathrm{e}^{2x-y}$.
  \item Classify $\Hess = \begin{pmatrix}2&4\\4&2\end{pmatrix}$.
  \item Find and classify the critical points of
        $f(x,y) = x^{2} - y^{2} + 6x - 4y$.
  \item Compute all second partials of $f(x,y) = \mathrm{e}^{x^{2}+y}$ and
        verify Young's theorem.
  \item Construct the Hessian of the Cobb--Douglas function
        $f(K,L) = K^{1/3}L^{2/3}$, $K,L>0$.
  \item Classify $\Hess = \begin{pmatrix}7&-3\\-3&1\end{pmatrix}$.
  \item Find and classify the critical points of
        $f(x,y) = x^{3} + y^{3} - 12x - 27y$.
  \item Compute all second partials of $f(x,y) = 1/(x+y)$ and verify
        Young's theorem.
  \item Construct the Hessian of
        $f(x,y,z) = x^{2} + 2y^{2} + 3z^{2} - xy + 2yz$.
  \item Classify $\Hess = \begin{pmatrix}-3&1\\1&-3\end{pmatrix}$.
  \item Find and classify the critical points of
        $f(x,y) = x^{2} + 4xy + y^{2} + 2y$.
  \item Compute all second partials of $f(x,y) = x^{1/2}y^{1/2}$
        ($x,y>0$) and verify $f_{xy} = f_{yx}$.
  \item Construct the Hessian of $f(x,y) = \ln(x^{2}+y^{2})$ at a generic
        point $(x,y)\neq(0,0)$.
  \item Find and classify the critical points of
        $f(x,y) = x^{4} + y^{4} - 4xy$.
  \item Compute all second partials of $f(x,y) = x\mathrm{e}^{y} + y\mathrm{e}^{x}$
        and verify Young's theorem.
  \item Construct the Hessian of $f(x,y) = x^{3} - 3xy^{2}$ (a harmonic
        function).
  \item Classify $\Hess = \begin{pmatrix}5&0\\0&-2\end{pmatrix}$ and state
        what this implies if $H$ is the Hessian at a critical point.
  \item Find and classify the critical points of
        $f(x,y) = xy - x^{2} - y^{2} + 6x + 3y$.
  \item Compute all second partials of $f(x,y) = \ln(xy)$ ($x,y>0$).
  \item Construct the Hessian of $f(x,y) = (x^{2}+y^{2})^{2}$ at a generic
        point.
  \item Classify $\Hess = \begin{pmatrix}-1&2\\2&-1\end{pmatrix}$.
  \item Find and classify the critical points of
        $f(x,y) = x^{3} - 3xy + y^{3}$.
  \item Compute all second partials of $f(x,y) = \sin(x)\cos(y)$ and
        verify Young's theorem.
  \item Construct the Hessian of $f(x,y,z) = \mathrm{e}^{xyz}$ at the point
        $(1,1,1)$ (evaluate numerically).
  \item Classify $\Hess = \begin{pmatrix}4&-2\\-2&4\end{pmatrix}$.
  \item Find and classify the critical points of the profit function
        $\pi(L,K) = 10L^{1/2}K^{1/2} - 2L - 3K$ (assume $L,K>0$).
  \item Compute all second partials of $f(x,y) = x^{2}y^{3}$ and verify
        Young's theorem.
  \item Construct the Hessian of
        $f(x,y) = \mathrm{e}^{x}\sin(y) - \mathrm{e}^{x}\cos(y)$.
  \item Classify $\Hess = \begin{pmatrix}9&6\\6&4\end{pmatrix}$. (Note the
        boundary case.)
  \item Find and classify the critical points of
        $f(x,y) = x^{2} + xy + y^{2} - 3x$.
  \item Compute all second partials of $f(x,y) = \arctan(y/x)$ ($x>0$) and
        verify Young's theorem.
  \item Construct the Hessian of the CES function
        $f(K,L) = (K^{1/2}+L^{1/2})^{2}$ at a generic point.
  \item Find and classify the critical points of
        $f(x,y) = -x^{4} - y^{4} + 8xy$.
  \item Compute all second partials of $f(x,y) = x^{2}\mathrm{e}^{-y}$ and
        verify Young's theorem.
  \item Construct the Hessian of $f(x,y) = \cos(xy)$ at the origin
        $(0,0)$.
  \item Classify $\Hess = \begin{pmatrix}0&3\\3&0\end{pmatrix}$.
  \item Find and classify the critical points of
        $f(x,y) = x^{3} + 3xy^{2} - 15x - 12y$.
  \item Construct the Hessian of $f(x,y) = \ln(1+x^{2}+y^{2})$ at the
        origin.
  \item Find and classify the critical points of
        $f(x,y) = \mathrm{e}^{-(x^{2}+y^{2})}$.
\end{problems}

\section{Solutions}

\begin{solutions}
  \item $f_x = 12x^{2} - 10xy$, $f_y = -5x^{2} + 28y^{3}$. Then
        $f_{xx} = 24x - 10y$, $f_{yy} = 84y^{2}$, and
        \[
          f_{xy} = \pd{}{y}\bigl(12x^{2}-10xy\bigr) = -10x, \qquad
          f_{yx} = \pd{}{x}\bigl(-5x^{2}+28y^{3}\bigr) = -10x . \; \checkmark
        \]

  \item $f_x = 10x + 6y - 3$, $f_y = 6x + 4y + 4$. Hessian
        \[
          \Hess = \begin{pmatrix}10&6\\6&4\end{pmatrix} .
        \]

  \item $D_1 = 6 > 0$, $D_2 = 18 - 4 = 14 > 0$. \textbf{PD.}

  \item $\grad f = (2x-6,\,8y+8) = (0,0) \Rightarrow x=3,\,y=-1$. Hessian
        $\Hess = \begin{pmatrix}2&0\\0&8\end{pmatrix}$: $D_1=2>0$,
        $D_2=16>0$, \textbf{PD}, so $(3,-1)$ is a strict \textbf{local
        minimum} with $f(3,-1) = 9+4-18-8+11 = -2$.

  \item Chain rule: $f_x = 3(x+2y)^{2}$, $f_y = 6(x+2y)^{2}$. So
        $f_{xx} = 6(x+2y)$, $f_{yy} = 24(x+2y)$, and
        \[
          f_{xy} = 12(x+2y) = f_{yx} . \; \checkmark
        \]

  \item $f_x = 3x^{2} + 2y^{2}$, $f_y = 4xy - 4$. Hessian
        \[
          \Hess = \begin{pmatrix}6x&4y\\4y&4x\end{pmatrix} .
        \]

  \item $D_1 = -4 < 0$, $D_2 = 20 - 4 = 16 > 0$. \textbf{ND.}

  \item $\grad f = (-2x+8,\,-6y-6) = (0,0) \Rightarrow x=4,\,y=-1$. Hessian
        $\Hess = \begin{pmatrix}-2&0\\0&-6\end{pmatrix}$: $D_1=-2<0$,
        $D_2=12>0$, \textbf{ND}, so $(4,-1)$ is a strict \textbf{local
        maximum} with $f(4,-1) = -16-3+32+6+5 = 24$.

  \item $f_x = 2x\ln(y)$, $f_y = x^{2}/y$. So $f_{xx} = 2\ln(y)$,
        $f_{yy} = -x^{2}/y^{2}$, and
        \[
          f_{xy} = 2x/y = f_{yx} . \; \checkmark
        \]

  \item $f_x = 2\mathrm{e}^{2x-y}$, $f_y = -\mathrm{e}^{2x-y}$. So
        $f_{xx} = 4\mathrm{e}^{2x-y}$, $f_{yy} = \mathrm{e}^{2x-y}$,
        $f_{xy} = -2\mathrm{e}^{2x-y}$. Hessian
        \[
          \Hess = \mathrm{e}^{2x-y}\begin{pmatrix}4&-2\\-2&1\end{pmatrix} .
        \]

  \item $D_1 = 2 > 0$, $D_2 = 4 - 16 = -12 < 0$. \textbf{Indefinite
        (saddle).}

  \item $\grad f = (2x+6,\,-2y-4) = (0,0) \Rightarrow x=-3,\,y=-2$. Hessian
        $\Hess = \begin{pmatrix}2&0\\0&-2\end{pmatrix}$: $D_2 = -4 < 0$,
        \textbf{indefinite}, so $(-3,-2)$ is a \textbf{saddle point}.

  \item Let $E = \mathrm{e}^{x^{2}+y}$. Then $f_x = 2xE$, $f_y = E$. So
        $f_{xx} = 2E + 4x^{2}E = (2+4x^{2})E$, $f_{yy} = E$, and
        \[
          f_{xy} = \pd{}{y}(2xE) = 2xE, \qquad
          f_{yx} = \pd{}{x}(E) = 2xE . \; \checkmark
        \]

  \item $f_K = \tfrac13K^{-2/3}L^{2/3}$, $f_L = \tfrac23K^{1/3}L^{-1/3}$.
        So
        \[
          f_{KK} = -\tfrac29K^{-5/3}L^{2/3}, \quad
          f_{LL} = -\tfrac29K^{1/3}L^{-4/3}, \quad
          f_{KL} = \tfrac29K^{-2/3}L^{-1/3} = f_{LK} . \; \checkmark
        \]
        Hessian
        \[
          \Hess = \tfrac29
          \begin{pmatrix}
            -K^{-5/3}L^{2/3} & K^{-2/3}L^{-1/3} \\
            K^{-2/3}L^{-1/3} & -K^{1/3}L^{-4/3}
          \end{pmatrix} .
        \]

  \item $D_1 = 7 > 0$, $D_2 = 7 - 9 = -2 < 0$. \textbf{Indefinite.}

  \item $\grad f = (3x^{2}-12,\,3y^{2}-27) = (0,0) \Rightarrow x^{2}=4,\,
        y^{2}=9$, giving $(\pm2,\pm3)$. Hessian
        $\Hess = \begin{pmatrix}6x&0\\0&6y\end{pmatrix}$, so $D_1=6x$,
        $D_2=36xy$.
        \begin{itemize}
          \item $(2,3)$: $D_1=12>0$, $D_2=216>0$, \textbf{PD}: local
                minimum, $f = 8+27-24-81 = -70$.
          \item $(2,-3)$: $D_2=-216<0$, \textbf{saddle}.
          \item $(-2,3)$: $D_2=-216<0$, \textbf{saddle}.
          \item $(-2,-3)$: $D_1=-12<0$, $D_2=216>0$, \textbf{ND}: local
                maximum, $f = -8-27+24+81 = 70$.
        \end{itemize}

  \item Let $S=x+y$, $f=S^{-1}$, $f_x = -S^{-2} = f_y$. So
        $f_{xx} = 2S^{-3} = f_{yy}$, and
        \[
          f_{xy} = 2S^{-3} = f_{yx} = \frac{2}{(x+y)^{3}} . \; \checkmark
        \]

  \item $f_x = 2x-y$, $f_y = 4y-x+2z$, $f_z = 6z+2y$. Hessian
        \[
          \Hess = \begin{pmatrix}2&-1&0\\-1&4&2\\0&2&6\end{pmatrix} .
        \]

  \item $D_1 = -3 < 0$, $D_2 = 9 - 1 = 8 > 0$. \textbf{ND.}

  \item $\grad f = (2x+4y,\,4x+2y+2) = (0,0)$. From the first, $x=-2y$;
        substituting, $-8y+2y+2=0 \Rightarrow y=1/3$, $x=-2/3$. Hessian
        $\Hess = \begin{pmatrix}2&4\\4&2\end{pmatrix}$: $D_1=2>0$,
        $D_2=4-16=-12<0$, \textbf{indefinite: saddle point} at
        $(-2/3,1/3)$.

  \item $f_x = \tfrac12x^{-1/2}y^{1/2}$, $f_y = \tfrac12x^{1/2}y^{-1/2}$.
        So
        \[
          f_{xx} = -\tfrac14x^{-3/2}y^{1/2}, \quad
          f_{yy} = -\tfrac14x^{1/2}y^{-3/2}, \quad
          f_{xy} = \tfrac14x^{-1/2}y^{-1/2} = f_{yx} . \; \checkmark
        \]

  \item Let $S=x^{2}+y^{2}$, $f=\ln S$, $f_x = 2x/S$, $f_y = 2y/S$. By the
        quotient rule:
        \[
          f_{xx} = \frac{2S-2x\cdot2x}{S^{2}} = \frac{2(y^{2}-x^{2})}{(x^{2}+y^{2})^{2}}, \quad
          f_{yy} = \frac{2(x^{2}-y^{2})}{(x^{2}+y^{2})^{2}}, \quad
          f_{xy} = \frac{-4xy}{(x^{2}+y^{2})^{2}} .
        \]
        Hessian
        \[
          \Hess = \frac{2}{(x^{2}+y^{2})^{2}}
          \begin{pmatrix}y^{2}-x^{2} & -2xy \\ -2xy & x^{2}-y^{2}\end{pmatrix} .
        \]
        Note $f_{xx}+f_{yy}=0$: $\ln(x^{2}+y^{2})$ is harmonic.

  \item $\grad f = (4x^{3}-4y,\,4y^{3}-4x) = (0,0) \Rightarrow y=x^{3}$
        and $x=y^{3}=x^{9}$. So $x(x^{8}-1)=0$, giving $x\in\{-1,0,1\}$.
        Critical points: $(0,0),(1,1),(-1,-1)$. Hessian
        $\Hess = \begin{pmatrix}12x^{2}&-4\\-4&12y^{2}\end{pmatrix}$.
        \begin{itemize}
          \item $(0,0)$: $\Hess = \begin{pmatrix}0&-4\\-4&0\end{pmatrix}$,
                $D_2=-16<0$, \textbf{saddle}.
          \item $(1,1)$: $\Hess = \begin{pmatrix}12&-4\\-4&12\end{pmatrix}$,
                $D_1=12>0$, $D_2=144-16=128>0$, \textbf{PD}: local
                minimum, $f=1+1-4=-2$.
          \item $(-1,-1)$: same $H$, \textbf{PD}: local minimum, $f=-2$.
        \end{itemize}

  \item $f_x = \mathrm{e}^{y}+y\mathrm{e}^{x}$, $f_y = x\mathrm{e}^{y}+\mathrm{e}^{x}$. So
        $f_{xx} = y\mathrm{e}^{x}$, $f_{yy} = x\mathrm{e}^{y}$, and
        \[
          f_{xy} = \mathrm{e}^{y}+\mathrm{e}^{x} = f_{yx} . \; \checkmark
        \]

  \item $f_x = 3x^{2}-3y^{2}$, $f_y = -6xy$. So $f_{xx}=6x$,
        $f_{yy}=-6x$, $f_{xy}=-6y=f_{yx}$. Hessian
        \[
          \Hess = \begin{pmatrix}6x&-6y\\-6y&-6x\end{pmatrix} .
        \]
        Note $f_{xx}+f_{yy}=0$: $f$ is harmonic.

  \item $D_1 = 5 > 0$, $D_2 = -10 < 0$. \textbf{Indefinite.} At a critical
        point this would mean a \textbf{saddle}.

  \item $\grad f = (y-2x+6,\,x-2y+3) = (0,0)$. So $2x-y=6$ and
        $x-2y=-3$. From the first, $y=2x-6$; substituting:
        $x-2(2x-6)=-3 \Rightarrow -3x+12=-3 \Rightarrow x=5$, $y=4$.
        Hessian $\Hess = \begin{pmatrix}-2&1\\1&-2\end{pmatrix}$:
        $D_1=-2<0$, $D_2=4-1=3>0$, \textbf{ND}, so $(5,4)$ is a
        \textbf{local maximum} with $f(5,4) = 20-25-16+30+12 = 21$.

  \item $\ln(xy) = \ln x + \ln y$. So $f_x = 1/x$, $f_y = 1/y$,
        $f_{xx} = -1/x^{2}$, $f_{yy} = -1/y^{2}$, $f_{xy} = 0 = f_{yx}$.
        \; $\checkmark$

  \item $f_x = 2(x^{2}+y^{2})\cdot2x = 4x(x^{2}+y^{2})$,
        $f_y = 4y(x^{2}+y^{2})$. Then
        \[
          f_{xx} = 4(x^{2}+y^{2})+4x\cdot2x = 4(3x^{2}+y^{2}), \quad
          f_{yy} = 4(x^{2}+3y^{2}), \quad
          f_{xy} = 4x\cdot2y = 8xy = f_{yx} .
        \]
        Hessian
        \[
          \Hess = 4\begin{pmatrix}3x^{2}+y^{2}&2xy\\2xy&x^{2}+3y^{2}\end{pmatrix} .
        \]

  \item $D_1 = -1 < 0$, $D_2 = 1 - 4 = -3 < 0$. Since $D_2<0$,
        \textbf{indefinite.}

  \item $\grad f = (3x^{2}-3y,\,-3x+3y^{2}) = (0,0) \Rightarrow y=x^{2}$
        and $x=y^{2}=x^{4}$, so $x(x^{3}-1)=0$, giving $x\in\{0,1\}$.
        Critical points: $(0,0)$ and $(1,1)$. Hessian
        $\Hess = \begin{pmatrix}6x&-3\\-3&6y\end{pmatrix}$.
        \begin{itemize}
          \item $(0,0)$: $\Hess = \begin{pmatrix}0&-3\\-3&0\end{pmatrix}$,
                $D_2=-9<0$, \textbf{saddle}.
          \item $(1,1)$: $\Hess = \begin{pmatrix}6&-3\\-3&6\end{pmatrix}$,
                $D_1=6>0$, $D_2=36-9=27>0$, \textbf{PD}: local minimum,
                $f=1-3+1=-1$.
        \end{itemize}

  \item $f_x = \cos(x)\cos(y)$, $f_y = -\sin(x)\sin(y)$. So
        $f_{xx} = -\sin(x)\cos(y)$, $f_{yy} = -\sin(x)\cos(y)$, and
        \[
          f_{xy} = -\cos(x)\sin(y) = f_{yx} . \; \checkmark
        \]

  \item Let $u=xyz$, $f=\mathrm{e}^{u}$. Then $f_x = yz\mathrm{e}^{u}$,
        $f_y = xz\mathrm{e}^{u}$, $f_z = xy\mathrm{e}^{u}$. Direct
        partials: $f_{xx} = (yz)^{2}\mathrm{e}^{u}$,
        $f_{yy} = (xz)^{2}\mathrm{e}^{u}$, $f_{zz} = (xy)^{2}\mathrm{e}^{u}$.
        Cross-partials (using the product rule on
        $f_x = yz\,\mathrm{e}^{xyz}$):
        \[
          f_{xy} = z\mathrm{e}^{xyz} + yz\cdot xz\mathrm{e}^{xyz}
                 = z\mathrm{e}^{xyz}(1+xyz) .
        \]
        Similarly $f_{xz} = y\mathrm{e}^{xyz}(1+xyz)$ and
        $f_{yz} = x\mathrm{e}^{xyz}(1+xyz)$. (Symmetry $f_{xy}=f_{yx}$
        etc.\ confirms Young's theorem.) At $(1,1,1)$: $u=1$,
        $\mathrm{e}^{u}=\mathrm{e}$, $1+xyz=2$, and $yz=xz=xy=1$, so
        \[
          \Hess(1,1,1) = \mathrm{e}\begin{pmatrix}1&2&2\\2&1&2\\2&2&1\end{pmatrix} .
        \]

  \item $D_1 = 4 > 0$, $D_2 = 16 - 4 = 12 > 0$. \textbf{PD.}

  \item $\pi_L = 5L^{-1/2}K^{1/2}-2$, $\pi_K = 5L^{1/2}K^{-1/2}-3$.
        Setting both to zero gives $5K^{1/2}/L^{1/2}=2$ and
        $5L^{1/2}/K^{1/2}=3$. Multiplying these two equations: $25=6$,
        which is a contradiction. So $\pi$ has \textbf{no interior
        critical point}. Why? The production function
        $Q=10L^{1/2}K^{1/2}$ is homogeneous of degree 1 (constant returns
        to scale), so along any ray $(\lambda L,\lambda K)$ revenue and
        cost both scale linearly: $\pi(\lambda L,\lambda K)=\lambda\pi(L,K)$.
        Hence profit is either identically zero, or unbounded above
        (depending on sign of $\pi$ along some direction). With
        $\max_{L/K}(10\sqrt{LK}-2L-3K)/K$: setting $t=L/K$ gives
        $\pi/K = 10\sqrt{t}-2t-3$, maximised at $t=25/4$ giving
        $\pi/K = 25-25/2-3 = 19/2 > 0$. So $\pi\to+\infty$ along the ray
        $L=(25/4)K$. The Hessian test does not apply because there is no
        critical point to classify. (Lesson: an inconsistent
        first-order system is itself a meaningful answer about the
        geometry of the problem.)

  \item $f_x = 2xy^{3}$, $f_y = 3x^{2}y^{2}$. So $f_{xx} = 2y^{3}$,
        $f_{yy} = 6x^{2}y$, and
        \[
          f_{xy} = 6xy^{2} = f_{yx} . \; \checkmark
        \]

  \item Let $g=\sin y-\cos y$, so $f=\mathrm{e}^{x}g$. Then
        $f_x = \mathrm{e}^{x}g$, $f_y = \mathrm{e}^{x}(\cos y+\sin y)$. So
        $f_{xx} = \mathrm{e}^{x}g = \mathrm{e}^{x}(\sin y-\cos y)$,
        $f_{yy} = \mathrm{e}^{x}(-\sin y+\cos y) = -\mathrm{e}^{x}g$, and
        \[
          f_{xy} = \mathrm{e}^{x}(\cos y+\sin y) = f_{yx} . \; \checkmark
        \]
        Hessian
        \[
          \Hess = \mathrm{e}^{x}
          \begin{pmatrix}
            \sin y-\cos y & \cos y+\sin y \\
            \cos y+\sin y & \cos y-\sin y
          \end{pmatrix} .
        \]

  \item $D_1 = 9 > 0$, $D_2 = 36 - 36 = 0$. \textbf{Inconclusive
        (positive semi-definite boundary case)}: Sylvester gives no
        verdict; we cannot rule out PSD with a zero eigenvalue. (In fact
        $H = (3,2)^{T}(3,2)$, so it is PSD with one zero eigenvalue.)

  \item $\grad f = (2x+y-3,\,x+2y) = (0,0)$. From the second: $x=-2y$;
        substituting: $-4y+y-3=0 \Rightarrow y=-1$, $x=2$. Critical point
        $(2,-1)$. Hessian $\Hess = \begin{pmatrix}2&1\\1&2\end{pmatrix}$:
        $D_1=2>0$, $D_2=4-1=3>0$, \textbf{PD}, so $(2,-1)$ is a strict
        \textbf{local minimum} with $f(2,-1) = 4-2+1-6 = -3$.

  \item With $u=y/x$, $f=\arctan u$: $f_x = -y/(x^{2}+y^{2})$,
        $f_y = x/(x^{2}+y^{2})$. Let $S=x^{2}+y^{2}$. Then
        \[
          f_{xx} = \frac{2xy}{S^{2}}, \qquad
          f_{yy} = \frac{-2xy}{S^{2}}, \qquad
          f_{xy} = \frac{-S+2y^{2}}{S^{2}} = \frac{y^{2}-x^{2}}{(x^{2}+y^{2})^{2}}
                  = f_{yx} . \; \checkmark
        \]
        (And $f_{xx}+f_{yy}=0$: $\arctan(y/x)$ is harmonic.)

  \item Let $g=K^{1/2}+L^{1/2}$. From sec06 sol.\ 23, $f_K = g/K^{1/2}$,
        $f_L = g/L^{1/2}$. Then
        \[
          f_{KK} = \frac{\tfrac12K^{-1/2}\cdot K^{1/2}-g\cdot\tfrac12K^{-1/2}}{K}
                  = \frac{\tfrac12-g/(2K^{1/2})}{K}
                  = \frac{1}{2K}-\frac{g}{2K^{3/2}} .
        \]
        More cleanly: $f_K = 1+(L/K)^{1/2}$, so
        $f_{KK} = -\tfrac12L^{1/2}K^{-3/2}$, and by symmetry
        $f_{LL} = -\tfrac12K^{1/2}L^{-3/2}$. Cross-partial:
        $f_{KL} = \tfrac12K^{-1/2}L^{-1/2} = f_{LK}$. \; $\checkmark$
        Hessian
        \[
          \Hess = \tfrac12
          \begin{pmatrix}
            -L^{1/2}K^{-3/2} & K^{-1/2}L^{-1/2} \\
            K^{-1/2}L^{-1/2} & -K^{1/2}L^{-3/2}
          \end{pmatrix} .
        \]

  \item $\grad f = (-4x^{3}+8y,\,-4y^{3}+8x) = (0,0) \Rightarrow x^{3}=2y$
        and $y^{3}=2x$. Substituting: $y=x^{3}/2$, then
        $(x^{3}/2)^{3}=2x \Rightarrow x^{9}/8=2x \Rightarrow x(x^{8}-16)=0$,
        so $x\in\{0,\pm16^{1/8}\}=\{0,\pm\sqrt2\}$. Points: $(0,0)$,
        $(\sqrt2,\sqrt2)$ (since $y=x^{3}/2=2\sqrt2/2=\sqrt2$), and
        $(-\sqrt2,-\sqrt2)$. Hessian
        $\Hess = \begin{pmatrix}-12x^{2}&8\\8&-12y^{2}\end{pmatrix}$.
        \begin{itemize}
          \item $(0,0)$: $\Hess = \begin{pmatrix}0&8\\8&0\end{pmatrix}$,
                $D_2=-64<0$, \textbf{saddle}.
          \item $(\sqrt2,\sqrt2)$: $x^{2}=y^{2}=2$,
                $\Hess = \begin{pmatrix}-24&8\\8&-24\end{pmatrix}$,
                $D_1=-24<0$, $D_2=576-64=512>0$, \textbf{ND}: local
                maximum, $f=-4-4+16=8$.
          \item $(-\sqrt2,-\sqrt2)$: same $H$, $f=-4-4+16=8$. \textbf{Local
                maximum.}
        \end{itemize}

  \item $f_x = 2x\mathrm{e}^{-y}$, $f_y = -x^{2}\mathrm{e}^{-y}$. So
        $f_{xx} = 2\mathrm{e}^{-y}$, $f_{yy} = x^{2}\mathrm{e}^{-y}$, and
        \[
          f_{xy} = -2x\mathrm{e}^{-y} = f_{yx} . \; \checkmark
        \]

  \item $f_x = -y\sin(xy)$, $f_y = -x\sin(xy)$. Then
        $f_{xx} = -y^{2}\cos(xy)$, $f_{yy} = -x^{2}\cos(xy)$, and
        (product rule)
        \[
          f_{xy} = -\sin(xy) - y\cdot x\cos(xy) = -\sin(xy) - xy\cos(xy) .
        \]
        At $(0,0)$: $\sin 0=0$, $\cos 0=1$. So $f_{xx}(0,0)=0$,
        $f_{yy}(0,0)=0$, $f_{xy}(0,0)=0$. Hessian
        \[
          \Hess(0,0) = \begin{pmatrix}0&0\\0&0\end{pmatrix} .
        \]

  \item $D_1 = 0$. Sylvester's stated criterion is inconclusive because
        $D_1$ is not strictly positive or negative. Directly:
        $D_2 = 0-9 = -9 < 0$, which by the third clause of Sylvester
        (sufficient on its own) gives \textbf{indefinite.}

  \item $\grad f = (3x^{2}+3y^{2}-15,\,6xy-12) = (0,0)$. From the second:
        $xy=2$. From the first: $x^{2}+y^{2}=5$. Combined with $xy=2$:
        $(x+y)^{2}=5+4=9$ so $x+y=\pm3$; $(x-y)^{2}=5-4=1$ so
        $x-y=\pm1$. Solutions $(x,y)\in\{(2,1),(1,2),(-2,-1),(-1,-2)\}$.
        Hessian $\Hess = \begin{pmatrix}6x&6y\\6y&6x\end{pmatrix}$, so
        $D_1=6x$, $D_2=36(x^{2}-y^{2})$.
        \begin{itemize}
          \item $(2,1)$: $D_1=12>0$, $D_2=36(4-1)=108>0$, \textbf{PD}:
                local minimum, $f=8+6-30-12=-28$.
          \item $(1,2)$: $D_1=6>0$, $D_2=36(1-4)=-108<0$, \textbf{saddle}.
          \item $(-2,-1)$: $D_1=-12<0$, $D_2=108>0$, \textbf{ND}: local
                maximum, $f=-8-6+30+12=28$.
          \item $(-1,-2)$: $D_1=-6<0$, $D_2=-108<0$, $D_2<0$ gives
                \textbf{saddle}.
        \end{itemize}

  \item Let $S=1+x^{2}+y^{2}$, $f=\ln S$, $f_x = 2x/S$, $f_y = 2y/S$.
        Then (using quotient rule):
        \[
          f_{xx} = \frac{2S-2x\cdot2x}{S^{2}} = \frac{2(1+y^{2}-x^{2})}{S^{2}}, \quad
          f_{yy} = \frac{2(1+x^{2}-y^{2})}{S^{2}}, \quad
          f_{xy} = \frac{-4xy}{S^{2}} = f_{yx} . \; \checkmark
        \]
        At the origin $(0,0)$: $S=1$, so
        \[
          \Hess(0,0) = \begin{pmatrix}2&0\\0&2\end{pmatrix} .
        \]

  \item $\grad f = \mathrm{e}^{-(x^{2}+y^{2})}(-2x,-2y) = (0,0)
        \Rightarrow x=y=0$, since $\mathrm{e}^{-r^{2}}>0$. Critical point
        $(0,0)$. Compute the Hessian: $f_x = -2x\mathrm{e}^{-(x^{2}+y^{2})}$,
        and at $(0,0)$, $f_{xx} = -2\mathrm{e}^{0}+(-2x)(-2x)\mathrm{e}^{0}\big|_{0,0}=-2$.
        Similarly $f_{yy}(0,0)=-2$, and
        $f_{xy}(0,0) = (-2x)(-2y)\mathrm{e}^{0}\big|_{0,0}=0$. Hessian
        $\Hess(0,0) = \begin{pmatrix}-2&0\\0&-2\end{pmatrix}$: $D_1=-2<0$,
        $D_2=4>0$, \textbf{ND: strict local maximum} at $(0,0)$ with
        $f(0,0)=1$. (Indeed, this is the global maximum, since $f\to0$
        at infinity and $f>0$ everywhere.)
\end{solutions}

\section{Additional Problems}

\begin{problems}
  \item Compute all second partials of $f(x,y) = 3x^{3} - 4x^{2}y + 5y^{3}$
        and verify Young's theorem.
  \item Construct the Hessian of
        $f(x,y) = 3x^{2} - 4xy + 2y^{2} + x - 5y$.
  \item Classify the definiteness of
        $\Hess = \begin{pmatrix}5&-1\\-1&4\end{pmatrix}$ using Sylvester's
        criterion.
  \item Find and classify the critical points of
        $f(x,y) = 2x^{2} + 3y^{2} - 4x + 6y + 1$.
  \item Compute all second partials of $f(x,y) = (2x+y)^{4}$ and verify
        $f_{xy} = f_{yx}$.
  \item Construct the Hessian of $f(x,y) = \mathrm{e}^{x+2y}$.
  \item Classify $\Hess = \begin{pmatrix}-3&1\\1&-2\end{pmatrix}$.
  \item Find and classify the critical points of
        $f(x,y) = -2x^{2} - 4y^{2} + 8x + 8y - 3$.
  \item Compute all second partials of $f(x,y) = xy\ln y$ ($y>0$) and
        verify Young's theorem.
  \item Construct the Hessian of $f(x,y) = x^{3}y + xy^{3}$.
  \item Classify $\Hess = \begin{pmatrix}2&3\\3&4\end{pmatrix}$.
  \item Find and classify the critical points of
        $f(x,y) = x^{2} - 2xy + 3y^{2} + 2x - 6y$.
  \item Compute all second partials of the Cobb--Douglas
        $f(K,L) = K^{1/2}L^{1/2}$ ($K,L>0$) and verify Young's theorem.
  \item Find and classify the critical points of
        $f(x,y) = x^{3} + y^{3} - 3x - 3y$.
  \item Classify $\Hess = \begin{pmatrix}-5&0\\0&-3\end{pmatrix}$ and state
        what it implies if $H$ is the Hessian at a critical point.
\end{problems}

\section{Additional Solutions}

\begin{solutions}
  \item $f_x = 9x^{2} - 8xy$, $f_y = -4x^{2} + 15y^{2}$. Then
        $f_{xx} = 18x - 8y$, $f_{yy} = 30y$, and
        \[
          f_{xy} = -8x = f_{yx} . \; \checkmark
        \]

  \item $f_x = 6x - 4y + 1$, $f_y = -4x + 4y - 5$. Hessian
        \[
          \Hess = \begin{pmatrix}6&-4\\-4&4\end{pmatrix} .
        \]

  \item $D_1 = 5 > 0$, $D_2 = 20 - 1 = 19 > 0$. \textbf{PD.}

  \item $\grad f = (4x-4,\,6y+6) = (0,0) \Rightarrow x=1,\,y=-1$. Hessian
        $\Hess = \begin{pmatrix}4&0\\0&6\end{pmatrix}$: $D_1=4>0$,
        $D_2=24>0$, \textbf{PD}, so $(1,-1)$ is a strict \textbf{local
        minimum} with $f(1,-1) = 2+3-4-6+1 = -4$.

  \item Chain rule: $f_x = 4(2x+y)^{3}\cdot2 = 8(2x+y)^{3}$,
        $f_y = 4(2x+y)^{3}$. So $f_{xx} = 48(2x+y)^{2}$,
        $f_{yy} = 12(2x+y)^{2}$, and
        \[
          f_{xy} = 24(2x+y)^{2} = f_{yx} . \; \checkmark
        \]

  \item $f_x = \mathrm{e}^{x+2y}$, $f_y = 2\mathrm{e}^{x+2y}$. So
        $f_{xx} = \mathrm{e}^{x+2y}$, $f_{yy} = 4\mathrm{e}^{x+2y}$,
        $f_{xy} = 2\mathrm{e}^{x+2y}$. Hessian
        \[
          \Hess = \mathrm{e}^{x+2y}\begin{pmatrix}1&2\\2&4\end{pmatrix} .
        \]
        (Note $\det\Hess = 0$: this is rank-1 PSD, consistent with $f$
        being a function of the single combination $x+2y$.)

  \item $D_1 = -3 < 0$, $D_2 = 6 - 1 = 5 > 0$. \textbf{ND.}

  \item $\grad f = (-4x+8,\,-8y+8) = (0,0) \Rightarrow x=2,\,y=1$. Hessian
        $\Hess = \begin{pmatrix}-4&0\\0&-8\end{pmatrix}$: $D_1=-4<0$,
        $D_2=32>0$, \textbf{ND}, so $(2,1)$ is a strict \textbf{local
        maximum} with $f(2,1) = -8-4+16+8-3 = 9$.

  \item $f_x = y\ln y$, $f_y = x\ln y + x = x(\ln y + 1)$. So
        $f_{xx} = 0$, $f_{yy} = x/y$, and
        \[
          f_{xy} = \ln y + 1 = f_{yx} . \; \checkmark
        \]

  \item $f_x = 3x^{2}y + y^{3}$, $f_y = x^{3} + 3xy^{2}$. So
        $f_{xx} = 6xy$, $f_{yy} = 6xy$,
        $f_{xy} = 3x^{2} + 3y^{2} = f_{yx}$. Hessian
        \[
          \Hess = \begin{pmatrix}
            6xy & 3x^{2}+3y^{2} \\
            3x^{2}+3y^{2} & 6xy
          \end{pmatrix} .
        \]

  \item $D_1 = 2 > 0$, $D_2 = 8 - 9 = -1 < 0$. \textbf{Indefinite
        (saddle).}

  \item $\grad f = (2x-2y+2,\,-2x+6y-6) = (0,0)$. From the first:
        $x = y-1$. Substitute: $-2(y-1)+6y-6 = 4y-4 = 0 \Rightarrow y=1$,
        $x=0$. Hessian $\Hess = \begin{pmatrix}2&-2\\-2&6\end{pmatrix}$:
        $D_1=2>0$, $D_2=12-4=8>0$, \textbf{PD}, so $(0,1)$ is a strict
        \textbf{local minimum} with $f(0,1) = 0-0+3+0-6 = -3$.

  \item $f_K = \tfrac12K^{-1/2}L^{1/2}$, $f_L = \tfrac12K^{1/2}L^{-1/2}$.
        So
        \[
          f_{KK} = -\tfrac14K^{-3/2}L^{1/2}, \quad
          f_{LL} = -\tfrac14K^{1/2}L^{-3/2}, \quad
          f_{KL} = \tfrac14K^{-1/2}L^{-1/2} = f_{LK} . \; \checkmark
        \]

  \item $\grad f = (3x^{2}-3,\,3y^{2}-3) = (0,0) \Rightarrow x^{2}=1,\,
        y^{2}=1$, giving $(\pm1,\pm1)$. Hessian
        $\Hess = \begin{pmatrix}6x&0\\0&6y\end{pmatrix}$, so $D_1=6x$,
        $D_2=36xy$.
        \begin{itemize}
          \item $(1,1)$: $D_1=6>0$, $D_2=36>0$, \textbf{PD}: local
                minimum, $f = 1+1-3-3 = -4$.
          \item $(1,-1)$: $D_2=-36<0$, \textbf{saddle}.
          \item $(-1,1)$: $D_2=-36<0$, \textbf{saddle}.
          \item $(-1,-1)$: $D_1=-6<0$, $D_2=36>0$, \textbf{ND}: local
                maximum, $f = -1-1+3+3 = 4$.
        \end{itemize}

  \item $D_1 = -5 < 0$, $D_2 = 15 > 0$. \textbf{ND}: at a critical point
        this would mean a \textbf{strict local maximum.}
\end{solutions}
